\section{Fase, registro dodecafásico y realización excepcional del estado común}
\label{sec:testigos-descendentes-ch}

El teorema de la sección~\ref{sec:cierre-cardinal-nativo} decide la cuestión
cardinal dentro del estado HMT completo. Las construcciones que siguen no son
un apéndice ontológico separable: hacen explícitas coordenadas correlacionadas
del mismo estado enriquecido. La doble proyección conserva informaciones
complementarias: valor arquimediano en una rama e incidencia, frontera y
genealogía en la otra.

La holonomía \(\Gamma_9\) compone nueve actualizaciones \(U_t\) sobre el estado
enriquecido. El núcleo reducido \(\mathcal K_{\mathrm{ph}}\) es una representación
posterior sobre emisiones ya construidas que conserva fase y memoria reducidas:
\[
 U_t\longrightarrow\Gamma_9\longrightarrow\mathcal K_{\mathrm{ph}},
 \qquad
 g(k+9)=g(k),\quad
 q_{\mathrm{mem},k+9}=q_{\mathrm{mem},k}+1,\quad
 \widehat x_{k+9}\neq\widehat x_k.
\]
La fase retorna, la memoria avanza y el estado no se reinicia. En cambio, el
sello dodecafásico entero se construye mediante el lector firmado y la
transformada de Hadamard:
\[
 \mathcal S_{12}\;(=R_{\mathrm{sgn}}):
 \mathcal X_{\mathrm{TPK}}^{\mathrm{term,enr}}
 \longrightarrow
 \mathcal U_{12}^{\mathrm{sgn}}:=\operatorname{im}\mathcal S_{12},
 \qquad
 U_{\mathrm{term}}\in\mathcal U_{12}^{\mathrm{sgn}},
 \qquad
 K:=\frac14\mathcal H_{12}U\in\mathbb Z^{12}.
\]
La coordenada \(U_{\mathrm{term}}\) se obtiene del estado enriquecido mediante
el registro orientado y los lectores de canal:
\[
 \mathcal S_{12}(x_{\mathrm{term}})
 =\Pi_H\!\left(
   \mathcal E_{108}^{90,120}
   (\mathcal R_{12}(x_{\mathrm{term}}))\right)
 =U_{\mathrm{term}}.
\]
La sección~\ref{sec:k-registro} desarrolla esa evaluación, la obtención
reversible de \(K\) y el retorno transversal a los mismos veinticinco canales.
La serialización de \(C_{\mathrm{term}}\) que consumen los controles Python es
un punto de entrada de comprobación, no un generador matemático ni una
restricción de la prueba impresa.
El vector \(K\) no produce por sí solo ninguna de las dos ramas. Junto con
las representaciones correlacionadas \(P,E,\Phi\), alimenta el cierre entero
\[
 \begin{gathered}
 Z_m:=P_m+E_m-\Phi_m-K_m,\qquad
 R_{N+1}(Z):=\sum_{r\geq0}Z_{N+1+r}1000^{-r-1},\\
 c_{N+1}^{\infty}:=\left\lfloor R_{N+1}(Z)\right\rfloor,\qquad
 s_m=Z_m+c_{m+1}^{\infty},\qquad
 A_m=s_m\bmod1000,\qquad
 c_m=\frac{s_m-A_m}{1000},\\
 \bigl(P_{1:N},E_{1:N},\Phi_{1:N},K_{1:N};
       c_{N+1}=c_{N+1}^{\infty}\bigr)_{N\geq1}
 \xrightarrow{\ \mathsf T_{\alpha,\infty}\ }
 (\alpha_N)_N
 \longmapsto
 \alpha_{\mathrm{HMT}}:=\varprojlim_N\alpha_N.
 \end{gathered}
\]
La condición \(c_{13}=0\) pertenece exclusivamente a la ejecución aislada de
doce bloques y expresa el racional \(\alpha_{12}^{(0)}\). No es la frontera
de la flecha infinita: en ésta cada horizonte recibe el cociente determinado
por su cola firmada; en particular, la primera vuelta de la sección completa
tiene \(c_{13}^{\infty}=-1\).
La representación excepcional conserva un codominio distinto:
\[
 \begin{gathered}
 B_K=\beta_{\mathrm{arq}}(K)=\{i:K_i\ge729\},\qquad
 s^{(0)}=P+E-\Phi-K,\\
 H^-_\alpha=\eta_{\mathrm{arq}}(s^{(0)})
 =\eta_{\mathrm{exc}}(B_K,P_\pi,H_e,H_\varphi,H_{\mathrm{APP}}),\\
 \widetilde{\mathcal I}_*
 =(B_K\subset H^-_\alpha;P_\pi,H_e,H_\varphi,H_{\mathrm{APP}})
 \longmapsto o_*=1
 \longmapsto v_\alpha=4\rho_1+\rho_2+\cdots+\rho_{12},\\
 N_{\alpha,0}
 =\{x\in N(C_W):\langle x,v_\alpha\rangle\equiv0\pmod3\},\qquad
 N_\alpha=N_{\alpha,0}+\mathbb Z\frac{v_\alpha}{3}\cong\Lambda_{24},\\
 C_W\longmapsto
 N(C_W):=\bigcup_{c\in C_W}\bigl(A_2^{12}+\phi(c)\bigr)
 \xrightarrow{\ \operatorname{Nbr}_{v_\alpha}\ }
 N_\alpha\cong\Lambda_{24}
 \longrightarrow V^\natural
 \xrightarrow{\ \operatorname{Aut}\ }\mathbb M.
 \end{gathered}
\]
El escalar \(\alpha_{\mathrm{HMT}}\), la cara \(B_K\), la hexada
\(H^-_\alpha\), el vector \(v_\alpha\) y el retículo \(N_\alpha\) pertenecen
a codominios diferentes. Comparten procedencia; no son objetos intercambiables.

\begin{proposition}[Unidad causal de las representaciones correlacionadas]
\label{prop:posterioridad-testigos-ch}
Los objetos
\[
 \mathcal K_{\mathrm{ph}},\quad K,\quad
 \pi_{\mathrm{HMT}},\varphi_{\mathrm{HMT}},e_{\mathrm{HMT}},
 \alpha_{\mathrm{HMT}},\quad H^-_\alpha,\quad v_\alpha,\quad
 N_\alpha,\quad V^\natural
\]
no son datos externos ni selectores retroactivos de la inferencia cardinal.
Todos se producen como coordenadas, representaciones o realizaciones de las dos
proyecciones del único estado enriquecido. La comparación cardinal no consulta
sus valores convencionales como oráculos, pero sí opera sobre la estructura
HMT completa que conserva su genealogía, sus regiones y su memoria.
\end{proposition}

\begin{proof}
La demostración cardinal utiliza el sistema inverso de historias, la recta
integral, la clausura semántica, el clasificador ordinal y las dos inyecciones
dentro del mismo estado APP--TRIT--TPK. Ninguna cifra convencional determina
esas aplicaciones. En el orden constructivo interno,
\(\mathcal K_{\mathrm{ph}}\) factoriza la memoria de \(\Gamma_9\); el sello
\(K\) se obtiene reversiblemente desde la coordenada terminal producida por
\(\mathcal E_{108}^{90,120}\circ\mathcal R_{12}\); la
cuatrirrelación se expresa desde las secciones coinductivas; y la bandera, el
vector \(v_\alpha\), el vecino de Leech y \(V^\natural\) se obtienen por la
proyección excepcional. Las secciones siguientes proporcionan las
construcciones completas y sus pruebas.
\end{proof}

Se preservan, en particular, las separaciones de tipo
\[
 \mathcal K_{\mathrm{ph}}\neq K,\qquad
 \mathcal H_{12}\neq H^-_\alpha,\qquad
 \alpha_{\mathrm{HMT}}\neq v_\alpha\neq N_\alpha.
\]
Tampoco existe una flecha canónica \(W_{24}\to\Lambda_{24}\): el diseño de
Witt y el vecino ternario de Leech son prolongaciones distintas del código
común. El grupo Monstruo actúa sobre \(V^\natural\), no sobre cifras decimales,
rutas TPK ni códigos de \(\alpha\).
